\begin{derivation}[从\eqnrefs{eq:sm-gauge-lagrangian-dp11}到\eqnrefs{eq:sm-wwv-vertex-dp11}]
一般 Yang--Mills 三点张量由场强的非线性项固定。对 $SU(2)_L$，生成元满足

\[
 [t^a,t^b]=\ii\epsilon^{abc}t^c,
\]
所以结构常数就是 $f^{abc}=\epsilon^{abc}$。由于 $\epsilon^{abc}$ 在任意两个指标相同时为零，含一对 $W^+,W^-$ 的非零三点项必须由原始 $\mathcal W^1,\mathcal W^2$ 分量组成，而第三条原始中性规范场是 $\mathcal W^3$。把带电场定义的逆变换
\begin{equation*}
 \mathcal W^1=\frac{\mathcal W^++\mathcal W^-}{\sqrt2},
 \qquad
 \mathcal W^2=\frac{\ii(\mathcal W^+-\mathcal W^-)}{\sqrt2}
\end{equation*}
代入群因子。固定第一条外腿为 $W^+$、第二条为 $W^-$、第三条为 $\mathcal W^3$ 时，两个非零实基分量给
\begin{align*}
 \sum_{a,b=1}^{2}C_a^+C_b^-\epsilon^{ab3}
 &=\frac1{\sqrt2}\left(-\frac{\ii}{\sqrt2}\right)\epsilon^{123}
 +\frac{\ii}{\sqrt2}\frac1{\sqrt2}\epsilon^{213}\\
 &=-\frac{\ii}{2}-\frac{\ii}{2}\\
 &=-\ii.
\end{align*}
一般实基顶角的群因子是 $-g\epsilon^{ab3}$，所以带电基中 $W^+W^-\mathcal W^3$ 的总体因子变为 $+\ii g$。洛伦兹张量仍保持
\begin{equation*}
 \mathcal V_{\mu\nu\rho}(p,q,r)
 =\eta_{\mu\nu}(p-q)_\rho
 +\eta_{\nu\rho}(q-r)_\mu
 +\eta_{\rho\mu}(r-p)_\nu.
\end{equation*}
再用
\begin{equation*}
 \mathcal W^3=s_WA+c_W\mathcal Z
\end{equation*}
得到
\begin{equation*}
 g\mathcal W^3=gs_WA+gc_W\mathcal Z.
\end{equation*}
由\eqnrefs{eq:gem-definition-dp11} 的 $gs_W=g_{\rm em}$ 代入上式，得到
\begin{equation*}
 g_{WW\gamma}=g_{\rm em},
 \qquad
 g_{WWZ}=gc_W,
\end{equation*}
即\eqnrefs{eq:sm-wwv-couplings-dp11}。把同一 $\mathcal V_{\mu\nu\rho}$ 的三条动量依次指定为 $p_+,p_-,k$，并保持全部流入约定，便得到\eqnrefs{eq:sm-wwv-vertex-dp11}；其中显式的 $+\ii$ 正来自上述实基到带电基的复线性变换。
\end{derivation}

\begin{derivation}[从\eqnrefs{eq:sm-wwv-vertex-dp11}到\eqnrefs{eq:sm-wwgamma-ward-dp41}]
正文\eqnrefs{eq:sm-wwv-vertex-dp11} 已固定三矢量顶角的动量结构。取 $V=\gamma$ 并用光子四动量收缩，再借助三条外腿全部流入的动量守恒，就能把结果整理成两条 $W$ 逆传播子之差，即正文\eqnrefs{eq:sm-wwgamma-ward-dp41}。

从\eqnrefs{eq:sm-wwv-vertex-dp11} 取 $V=\gamma$ 并与光子动量 $k^\rho$ 收缩：
\begin{align*}
 \frac{k^\rho\Gamma^{WW\gamma}_{\mu\nu\rho}}{+\ii g_{\rm em}}
 ={}&\eta_{\mu\nu}k\cdot(p_+-p_-)
 +k_\nu(p_--k)_\mu
 +k_\mu(k-p_+)_\nu.
\end{align*}
动量守恒 $k=-p_+-p_-$ 给出第一项
\begin{align*}
 k\cdot(p_+-p_-)
 &=(-p_+-p_-)\cdot(p_+-p_-)\\
 &=-p_+^2+p_-^2.
\end{align*}
后两项分别展开为
\begin{align*}
 k_\nu(p_--k)_\mu
 &=-(p_{+\nu}+p_{-\nu})(p_{+\mu}+2p_{-\mu}),\\
 k_\mu(k-p_+)_\nu
 &=(p_{+\mu}+p_{-\mu})(2p_{+\nu}+p_{-\nu}).
\end{align*}
把乘积逐项展开，$p_{+\mu}p_{-\nu}$ 与 $p_{-\mu}p_{+\nu}$ 的交叉项互相抵消，剩下
\begin{equation*}
 k_\nu(p_--k)_\mu+k_\mu(k-p_+)_\nu
 =p_{+\mu}p_{+\nu}-p_{-\mu}p_{-\nu}.
\end{equation*}
因此
\begin{align*}
 k^\rho\Gamma^{WW\gamma}_{\mu\nu\rho}
 =+\ii g_{\rm em}\Big[
 &(p_-^2\eta_{\mu\nu}-p_{-\mu}p_{-\nu})\\
 &-(p_+^2\eta_{\mu\nu}-p_{+\mu}p_{+\nu})
 \Big].
\end{align*}
在两个括号中同时减去相同的质量项 $m_W^2c^2\eta_{\mu\nu}$ 不改变差值，于是恰好得到
\begin{equation*}
 k^\rho\Gamma^{WW\gamma}_{\mu\nu\rho}
 =+\ii g_{\rm em}
 \left[\Pi_{\mu\nu}(p_-)-\Pi_{\mu\nu}(p_+)\right],
\end{equation*}
即\eqnrefs{eq:sm-wwgamma-ward-dp41}。若两条 $W$ 外腿在壳，则
\begin{equation*}
 p_\pm^2=m_W^2c^2,
 \qquad
 p_\pm\cdot\varepsilon_\pm=0,
\end{equation*}
所以 $\varepsilon_+^\mu\Pi_{\mu\nu}(p_+)=0$ 且 $\Pi_{\mu\nu}(p_-)\varepsilon_-^\nu=0$。这给出外部光子纵向替换后的零振幅，并同时检查了动量号、质量项与偏振条件。
\end{derivation}
